\documentclass[10pt]{article}
\usepackage[margin=20mm]{geometry}
\usepackage{amsmath,amssymb,mathtools,fontspec,unicode-math,xurl,hyperref}
\IfFontExistsTF{DejaVu Serif}{\setmainfont{DejaVu Serif}}{\setmainfont{Latin Modern Roman}}
\setmathfont{latinmodern-math.otf}
\setlength{\parindent}{0pt}\setlength{\parskip}{6pt}\sloppy
\title{Bailey Transformations, Character Formulas
and Proof Corrections}
\author{AI-assisted technical review}
\date{}
\begin{document}
\maketitle
Abstract: The exact 51-page author manuscript arXiv:2303.09069v3 is examined. Ten mathematical discrepancies are established by finite calculations, coefficient comparisons or direct parameter substitutions. One finding is assessed as Major: the printed proof of the central Bailey kernel identity passes through a false double-sum expansion, and the identity is subsequently used in the Bailey-tree argument. A corrected finite expansion repairs that derivation, without establishing that the printed step was valid. Nine further findings are Minor within their demonstrated scopes. The kernel identity and principal results are not disproved. Prior-theorem citations remain counterevidence, but their precise equivalence has not been verified from a retained complete source. Scientific adjudication remains pending.\par
\section*{Source, identity and review boundary}
The frozen record M25-0254 is S. Ole Warnaar, An \(\mathrm A_2\) Bailey tree and \(\mathrm A_2^{(1)}\) Rogers–Ramanujan-type identities, doi: 10.4171/JEMS/1627 (https://doi.org/10.4171/JEMS/1627). The actual PDF is arXiv:2303.09069v3, submitted on 24 February 2025 at 00:57:36 UTC; the first page also prints 24 February 2025 [source]. The publisher records submission on 19 April 2023, acceptance on 19 February 2025 and online publication on 14 April 2025 [record]. The author manuscript is not assumed identical to the publisher’s version. The frozen DOI and 2025 batch membership remain unchanged. Formula and page references below are to this exact manuscript, whose printed and PDF page numbers agree.\par
The initial review records reading all 51 pages, including the appendix and references. The present independent acceptance re-examined the reported Problems and their dependencies; it does not claim a second complete review of every page. The scope includes the conventions for negative factorial indices, the root-lattice support, the finite inverse matrices, the two- and four-parameter transformations, the Kanade–Russell iteration, the character formulas and the limiting arguments. Actual source images were checked for the evidence cited below. A bounded search of arXiv, the publisher and the author’s publication-page route did not locate a separate erratum, supplementary proof or required code. The publisher PDF route returned a subscription page and the author-page request was inaccessible. These are access and search limits, not evidence that no additional material exists. No author code was executed, and the absence of code is not assessed as a defect.\par
Throughout, \((a;q)_n\) and \(\theta(z;q)=(z,q/z;q)_\infty\) have the manuscript’s meaning. The arguments are primarily in formal series; analytic arguments below use \(0<|q|<1\) away from poles and then identify the corresponding formal coefficients.\par
\section*{Confirmed mathematical findings}
Problem 1 (Minor; p. 4, first unnumbered identity).  The prefactor for the \(\tau=-1\) specialization has the wrong modulus.\par
The printed prefactor is \((q^{2k+2};q^{2k+2})_\infty^2\), although the theta factors have modulus \(q^{3k+2}\). Take \(k=1\), \(a=1\), \(b=0\). These integers satisfy every restriction in that display, including \(b\ne k\) and the stated convention \(\mu_0=\infty\). Its left side reduces to \[L(q)=\sum_{n\ge0}\frac{q^{n^2}}{(q;q)_n(q;q)_\infty}.\] Indeed, \((1-q)(q^2;q)_\infty=(q;q)_\infty\). Up to degree four the inner series is \(1+q+q^2+q^3+2q^4\), since only \(n=0,1,2\) contribute. Multiplication by \((q;q)_\infty^{-1}=1+q+2q^2+3q^3+5q^4+O(q^5)\) gives \[L(q)=1+2q+4q^2+7q^3+13q^4+O(q^5).\] The printed right side, to the same order, is \[\frac1{(1-q)^2(1-q^2)(1-q^3)}+O(q^5)
 =1+2q+4q^2+7q^3+11q^4+O(q^5).\] Thus it fails by \(2q^4\). Replace \(2k+2\) by \(3k+2\) in the prefactor. Equation (1.5) and the correct product evaluation from (5.16) already use that modulus. The separate main identity and Theorem 1.2 are not refuted by this display error.\par
Problem 2 (Minor; p. 5, paragraph following Theorem 1.3).  The highest-weight label omits the level shift by three.\par
The paragraph assigns the series to \(L((K-a)\Lambda_0+a\Lambda_1)\). The manuscript’s own general assignment on p. 3 is \[(K-a-b-3)\Lambda_0+a\Lambda_1+b\Lambda_2.\] Setting \(b=0\) therefore gives \(L((K-a-3)\Lambda_0+a\Lambda_1)\). This is a numerical discrepancy in the representation label: each \(\Lambda_i\) has level one, so the printed weight has level \(K\) and the required weight has level \(K-3\). For the admissible example \(k=1,a=1,K=5\), the two labels are \(4\Lambda_0+\Lambda_1\) and \(\Lambda_0+\Lambda_1\); the central element acts by different scalars. There is also a coefficient discrepancy: (1.6a) gives \(\sum_{n\ge0}q^{n^2}/(q;q)_n
=1+q+q^2+O(q^3)\), whereas the printed level-five label, under the general product assignment on p. 3, requires modulus eight and gives \[\frac{(q^8;q^8)_\infty^2}{(q;q)_\infty^2}
 \theta(q^2,q,q^3;q^8)=1+q+2q^2+O(q^3).\] Thus this is not merely a harmless change of name for the same series. The correction is confined to this identification. The multisums (1.6a,b), the product expressions and the separately stated \(\mathcal W_3(3,K)\) conformal-weight formula are not changed. The correct general weight already printed on pp. 3 and 7 is direct counterevidence to interpreting this as a new level convention.\par
Problem 3 (Major; p. 17, proof of Theorem 4.2).  The double sum immediately below the transformed \({}_2\phi_1\) loses a factor and has an incorrect exponent.\par
The correct expansion is \[\frac1{(wq;q)_m}
 \sum_{s=0}^n\sum_{r=s}^n
 \frac{z^r w^s q^{r^2-rs+s^2}}
 {(q;q)_{n-r}(q;q)_{m-s}(zq;q)_r
  (q;q)_{r-s}(q,zwq;q)_s}.
 \tag{2.1}\] Terms with \(s>m\) are zero by the stipulated reciprocal-factorial convention. The printed display omits \((wq;q)_m^{-1}\) and instead has \(q^{\binom{r-s}{2}+r^2}\).\par
Both errors can be separated without approaching a pole. Put \(q=1/2\), \(z=2/3\), \(w=3/5\). At \(n=0,m=1\) the printed sum is \(2\), whereas the preceding expression and \(\Phi_{0,1}\) are \(20/7\). At \(n=2,m=0\) the missing prefactor is one, yet the printed sum is \(71/15\), whereas the preceding expression and \(\Phi_{2,0}\) are \(24/5\).\par
For a direct derivation of (2.1), use \[(q^{-r};q)_s=(-1)^s q^{-rs+\binom s2}
              \frac{(q;q)_r}{(q;q)_{r-s}}\] and the corresponding formula with \(m\). Multiplication by \((wq^{m+1})^s\) leaves the exponent \(r^2-rs+s^2\), and cancellation of \((q;q)_m\) leaves the displayed prefactor. After \(r\mapsto r+s\), this corrected expression gives the next \({}_1\phi_1\) line of the manuscript. This supplies a finite algebraic repair and supports the kernel identity, but the printed transition performs this shift on a false double sum. The supplied proof therefore contains a substantive failure in a central step. Equation (4.5) is then used on p. 19 for the constant, \(u\) and \(uv\) coefficients in the proof of Theorem 4.4, and for the shifted iteration underlying the Bailey tree. This actual dependency makes the finding Major; it does not make the kernel identity or the principal identities false.\par
The manuscript cites [8, Theorem 4.3], [89, Theorem 5.1] and [31, Corollary 7.9] as prior formulations on p. 16. These citations and the correct neighboring expressions are counterevidence to a false-theorem interpretation. No retained complete copy of those prior sources was available for this acceptance, so their precise equivalence has not been independently verified and is not counted as an established alternative proof. The corrected finite expansion above is the demonstrated repair.\par
Problem 4 (Minor; p. 30, first constant-term equality for \(A_{\nu;k}\)).  An extra column-index term makes the first constant-term expression false.\par
After using one consistent modulus symbol, the exponent should be \[rk\binom{y_i}{2}-kiy_i-(ry_i+j-i)\nu_j.
 \tag{2.2}\] The first constant-term line instead ends with \(+(ry_i+j-i)(j-\nu_j)\). Its next line already has the correct negative \(\nu_j\) term.\par
The text expressly generalizes the calculation to arbitrary rank. Take \(r=2\), \(k=2\) and \(\nu=(0,0)\). In the definition of \(A_{\nu;k}\) the two columns coincide, so every determinant is zero. In the erroneous constant-term expression, write \((y_1,y_2)=(m,-m)\). Direct expansion gives \[\sum_{m\in\mathbb Z}\bigl(q^{4m^2}-q^{(2m+1)^2}\bigr),\] whose constant coefficient is one, not zero. Removing the extra \(j\) gives (2.2), and multilinearity then yields exactly the following line. The subsequent transpose, change \(y\mapsto-y\) and Vandermonde evaluation are consistent; the extra linear term produced by negation vanishes because \(\sum_i y_i=0\). Also use \(k\), rather than a mixture of \(K\) and \(k\), in the definition of \(A_{\nu;k}\). This is an erroneous intermediate display, not a refutation of the recovered constant-term argument.\par
Problem 5 (Minor; p. 31, last display of the proof of Theorem 1.2).  The final product sign repeats an already expanded three-factor theta product.\par
The line contains \(\prod_{1\le i<j\le3}\theta(q^{a+1},q^{b+1},q^{a+b+2};q^K)\). The expression following the product sign is independent of \(i,j\); under the manuscript’s multiple-argument convention this cubes the required three-theta product. For \(k=1\), \(\tau=-1\), \(a=1\), \(b=0\) and \(K=5\), that three-theta product is \(1-q+O(q^2)\). The prefactor \((q^5;q^5)_\infty^2/(q;q)_\infty^3\) is \(1+3q+O(q^2)\). The intended result has coefficient \(2\) at \(q\); the printed result has coefficient \(0\). Delete the outer product sign. Alternatively retain the product with its unexpanded factor \(\theta(q^{\nu_i-\nu_j};q^K)\). The preceding general-rank formula and (1.5) both give the correct expression, so the endpoint of the proof is recovered directly.\par
Problem 6 (Minor; p. 32, parameter bounds preceding Theorem 6.1).  Two primed-parameter bounds are misstated.\par
Assume exactly the stated range \(2\le k<a\le\lfloor3k/2\rfloor\) and \(0\le b\le3k-2a\). For \[(a',b')=(2k-a,a+b-k),\qquad
 (a'',b'')=(2k-a-b-1,a-k-1),\] the text claims \(a'\le\lceil k/2\rceil\) and subsequently uses \(a'\) where \(a''\) is needed. The admissible choice \((k,a,b)=(4,5,0)\) gives \((a',b')=(3,1)\) and \((a'',b'')=(2,0)\). It contradicts both \(a'\le2\) and the second printed \(a'\le k-2=2\).\par
The bounds actually implied by the hypotheses are \[0<b'\le a'\le k-1,\qquad
 \lceil k/2\rceil\le a',\qquad
 0\le b''\le a''\le k-2.\] For example \(b'\le a'\) is just \(b\le3k-2a\), and \(a''-b''=3k-2a-b\ge0\). These correct bounds place both parameter pairs in the domain of Theorem 1.2. Consequently the application of the theta three-term relation and the statement of Theorem 6.1 survive.\par
Problem 7 (Minor; p. 32, first modulus-seven identity).  The factor \((q^7;q^7)_\infty\) needs exponent two.\par
The unnumbered identity for the below-the-line case \((a,b,k,\tau)=(2,0,1,1)\) has only one such factor. Writing its two summation indices as \(r,s\), separate the numerator into the two exponents \[E_1=r^2-rs+s^2-r+s,\qquad E_2=r^2-rs+s^2+r .\] The common reciprocal denominator is \((1-q)^{-1}(q;q)_r^{-1}(q;q)_s^{-1}(q^2;q)_{r+s}^{-1}\). Exact coefficient multiplication gives coefficient \(126\) at \(q^7\) on the left and \(127\) for the printed product. This calculation is finite: for coefficients through degree \(N\), \[E_1\ge\tfrac12((r-1)^2+(s+1)^2)-1,\qquad
 E_2\ge\tfrac12(r^2+s^2)+r.\] If \(r\) or \(s\) is at least \(N+3\), both exceed \(N\). The denominator introduces no negative powers. Thus the saved enumeration through \(N=10\), with \(0\le r,s\le12\), includes every contribution, including the cases \(2r-s<0\).\par
Squaring the prefactor gives coefficient \(126\) and agrees through degree ten in the exact computation. Its necessity also follows from the rank-two product normalization displayed on p. 31. The formula is a separately cited example; this correction does not alter the proof or domain of Theorem 6.1.\par
Problem 8 (Minor; p. 37, equation (7.6)).  The last index in the quadratic exponent must be \(j\).\par
The exponent is printed as \(\sum_{i,j=1}^6 n_iA_{ij}n_i\). The displayed matrix has row sums \[(11/2,\,11/2,\,9,\,10,\,12,\,14).\] Consequently the printed expression equals \(\sum_i s_i n_i^2\) with those row sums \(s_i\). Every nonzero tuple contributes only from degree at least \(11/2\), and the zero tuple contributes exactly one. Its coefficient of \(q\) is therefore zero. In contrast the defined character has product \[\widetilde\chi^8_{2,2}
 =\frac{(q^8;q^8)_\infty^2}{(q;q)_\infty^2}
   \theta(q^3,q^3,q^6;q^8)=1+2q+O(q^2).\] All parameters here are those of the displayed special case. Replace the final \(n_i\) by \(n_j\). Substitution of the six variables printed immediately above (7.6) into the preceding exponent yields precisely the displayed symmetric matrix with the corrected quadratic form; this was also checked by exact polynomial algebra. Thus the wrong printed series is not the announced character, but the matrix and its derivation support the intended identity. The general character theorems are not disproved.\par
Problem 9 (Minor; p. 42, equation (8.10b)).  The first argument of the subtracted character is \(z\), not \(q\).\par
Let \(F(z,w;q)=\operatorname{ch}W'_{\Lambda_0}(z,w;q)\). The printed left side is \(F(z,w;q)-F(q,wq;q)\). From (8.2) in the stated vacuum case \(k=1,a=b=0\), \[F(z,w;q)=\sum_{r,s\ge0}
 \frac{z^rw^sq^{r^2-rs+s^2}}{(q;q)_r(q;q)_s}.\] At the valid formal specialization \(z=w=0\), the printed left side is \[1-\sum_{r\ge0}\frac{q^{r^2+r}}{(q;q)_r}
 =-q^2+O(q^3),\] whereas the right side, having the factor \(wq\), is zero. The correction is \[F(z,w;q)-F(z,wq;q)=wqF(zq^{-1},wq^2;q).\] This follows term by term: multiply the summand by \(1-q^s\), cancel that factor against \((q;q)_s\), and put \(s=t+1\). The corrected recursion is also the \(z,w\) interchange of (8.10a). This discussion follows the completed proof of Theorem 8.1, so its typographical error does not invalidate that proof.\par
Problem 10 (Minor; p. 46, opening factorial identity).  The base of each factor on the left must be \(q^r\).\par
The printed equality is \(\prod_{i=1}^r(aq^{r-i};q)_k=(a;q)_{rk}\). Take \(r=k=2\), \(a=1/3\), \(q=1/2\), which are allowed nonzero parameters with no pole. The two sides are respectively \[\frac{275}{648}\quad\hbox{and}\quad\frac{1265}{2592},
 \qquad\hbox{with difference }-\frac{55}{864}.\] The correct identity is \[\prod_{i=1}^r(aq^{r-i};q^r)_k=(a;q)_{rk}.\] Its exponents \(r-i+rj\), for \(1\le i\le r\) and \(0\le j<k\), are exactly \(0,1,\ldots,rk-1\). The preceding instruction \(q\mapsto q^r\) supplies this base, and the following transformed summand uses the resulting contiguous products. The interpolation, rotation and coefficient extraction leading to (A.4) remain available with this correction; (A.4) itself was not refuted.\par
\section*{Supported proof structure and boundary checks}
The following observations distinguish the local errors from the central statements they accompany.\par
\subsection*{Finite kernels and the tree}
The proof of Lemma 3.3 is polynomial in \(u\). Its constant coefficient is the reproducing identity with \(z/q\); the linear coefficient follows by \(K_{n;r+1}(z/q)=zK_{n-1;r}(zq)\) and \(\Phi_{-1}=0\). The infinite-partition limit in Corollary 3.2 is coefficientwise: a nonzero terminal part after \(h\) iterations carries degree at least \(h\).\par
For the \(\mathrm A_2\) kernel, the shifted lower bounds in Corollaries 4.5, 4.10 and 4.11 were retained. The apparent ratios in (4.19) cancel into a six-term polynomial, so boundary values such as \(u=d=0\) are interpreted in that polynomial, not by dividing by zero. The two finite inverse-matrix products, the coefficient identities and the conversion of the specialized Gustafson summand were checked. The analytic bilateral identity is used first in its convergence region; its terminating specialization then gives a rational identity. These finite checks supplement the algebraic derivations and do not replace them.\par
\subsection*{Iteration to the main product}
Equation (5.7) uses \(t\ge2\); the manuscript explicitly identifies the exceptional support at \(t=1\), \(y=(-1,0,1)\). The proof uses \(t=3\), so that exception does not invalidate it. The stages use \(k-a\), \(a-b\) and \(b\) iterations in the range \(0\le b\le a\le k\). At \(a=b\) the central step is justified by (5.9), rather than by an unavailable instance with \(a-1\). Interchanging \(a,b\) is accompanied by \((n,m)\mapsto(m,n)\) and \(y\mapsto(-y_3,-y_2,-y_1)\).\par
The large-\(n,m\) limit can be justified uniformly. Expand (4.19) into its six terms and use (4.25), with \(Y=3y\). Each term contains a factorial ratio of the form \[\frac{(q;q)_{n+m}}
 {\prod_{i=1}^3
  (q;q)_{n-Y_i-\sigma_i+i}
  (q;q)_{m+Y_i+\sigma_i-i}},\] times a monomial whose exponent is linear in \(Y\). A negative reciprocal-factorial index gives zero; otherwise these ratios are uniformly bounded for fixed \(0<|q|<1\). The outer exponent has positive quadratic part \((3K/2)\sum_i y_i^2\), so Gaussian domination permits the limit. The multisum side has the positive quadratic form \(\lambda_i^2-\lambda_i\mu_i+\mu_i^2\). The corrected constant-term argument and Macdonald identity then give the stated product.\par
Separate exact checks tested all 39 choices with \(k=1,2\), all allowed \(a,b,\tau\), through \(q^{12}\), and 147 finite root identities at \(q=1/2\), without a mismatch in (1.5) or (5.2). No inference of universal correctness is drawn solely from those tests.\par
\subsection*{Characters and principal subspaces}
The \(\sigma_i\) in the p. 33 definition of \(v_i\) is unbound. Removing it gives \(v_i=\chi(i\ge b)-\lambda_i\) for \(i<k\), as follows directly from (1.5). These sequences satisfy the monotonicity hypotheses of Lemma 7.1 used in the character conversion. The endpoint \(\tau=-1,b=k\) is handled by the stated symmetry or the separate \(a=b=k\) expression, not by extending Theorem 7.3 beyond its displayed first range. Twenty-five instances of Theorems 7.2 and 7.3 at \(k=1,2\) were checked through \(q^{10}\). The quadratic-form substitution preceding (7.6) provides the specific counterevidence recorded in Problem 8.\par
For Theorem 8.1, the unit-pair identity (8.5) has finite support \(y_1\le n\), \(y_1+y_2\le m\) with \(y\in Q_+\). The proof follows the same three finite transformations and the endpoint \(a=b\) is explicitly supplied by (8.6). The manuscript itself defines the right side of (8.2) for \(a+b>k\) as an extension of notation; it does not silently assert that those weights are dominant.\par
For fixed \(a,b\), the nonzero-\(y\) terms in the bosonic formula have valuation at least \[(k-a-b+1)\sum_{y_i>0}y_i .\] For sufficiently large \(k\) this tends to infinity uniformly over nonzero \(y\in Q_+\), leaving the \(y=0\) determinant and proving the indicated coefficientwise passage to Corollary 1.5. On the partition side the positive quadratic tails likewise give finite support at each fixed degree. This addresses the limiting issue without imposing an unwritten analytic ordering on \(x_1,x_2,x_3\).\par
\subsection*{Appendix and general-rank discussion}
After the base correction in Problem 10, the specialized summand of (A.1) has support independent of \(N\). Clearing the indicated denominator gives polynomials in \(d\) of degree at most \(n\); their equality for \(d=q^{-rN}\) therefore proves the rational continuation used in the argument. The rotation sends a fixed total to the root lattice, and the finite \(q\)-binomial theorem gives (A.4). The root’s independently written calculation tested (A.4) at \(q=1/3\), \(r=2,3,4\), \(0\le m\le n\le4\), a total of 45 exact instances, all equal. Seven additional \(\mathrm A_1\) chain/tree checks at generic rational parameters also agree.\par
The reversal symmetry and chain (9.2)–(9.3) follow from the Cartan symmetry and splitting the first partition parts in (9.1). The general-rank tree and manifest positivity of the general root identity are presented as open problems. The review does not reinterpret them as proved assertions. The external Whittaker-vector and representation-theoretic results are dependencies whose entire original proofs were not independently reconstructed here.\par
\section*{Supporting checks}
Several remaining notational matters merit correction without turning them into manufactured theorem counterexamples. Conjecture 1.1 and Theorem 1.3 should specify \(k\ge1\) for their displayed terminal-index multisums, or supply a separate \(k=0\) definition. Merely assigning an arbitrary value to an undefined \(g_{\lambda_0,\mu_0}\) is not a valid disproof. On p. 38 the pairing uses the undefined printed symbol \(a_i^\vee\); if the intended symbol is \(\alpha_i^\vee\), the correct convention is \(\langle\Lambda_0,\alpha_i^\vee\rangle=\delta_{0i}\), as the next sentence and the level-one assertion require. The proof text on p. 19 also contains a common-factor label and an \(m\) for \(s\) index slip. The p. 26 inversion sentence should start from \(\tau=+1\) to derive \(\tau=-1\); the independently checked inversion exponents support that correction. The unbound \(\sigma_i\) on p. 33 is a notation issue, as explained above. These do not contradict the independently checked shifted identities.\par
Each of Problems 1–10 has a concrete mathematical discrepancy. Problem 3 is Major because the supplied proof of the central kernel identity contains a false expansion that is actually used by the subsequent transformations. Its explicit repair does not validate the printed proof. Problems 1–2 and 4–10 are Minor within the demonstrated scopes. In particular, for Problem 4 the explicit constant-term evaluation and multilinearity already yield the next correct line after deleting the erroneous intermediate display. For Problem 5 the correct general product and the stated \(r=3\), \(\nu=(a+b+2,b+1,0)\) substitution already give the required three-theta product. These are identifiable surviving arguments, rather than unproved assumptions supplied to conceal a core gap. The kernel identity and main results are not disproved. The exact equivalence of the cited prior kernel theorems remains unverified here. Scientific decision and adjudication remain pending; there are no human reviews. The current independent checks, their bounded native assistance and exact artifact bindings are recorded in the accompanying findings, coverage and provenance.\par
Their persistent names are historical; their current receipts identify M25-0254 and the current batch. Their role, source-page coverage, calculations and limitations are retained separately. They are not human referees.\par
That is attributed to the controller and does not independently attest the backend model. The Personal-account observation is likewise user-attributed to the controller, not personally observed by this reviewer. The session’s runtime self-description is recorded separately. The current build and actual report-page inspections are bound to the current source and report in the accompanying receipts; rendering alone is not inspection.\par
Exact source. bytes; SHA256\par
655e6748f5c59f7f4fa269ba7d107ed8400a0f71341b1a2b1e4797b6ee29dfc8\par
9 S. O. Warnaar, An \(\mathrm A_2\) Bailey tree and \(\mathrm A_2^{(1)}\) Rogers–Ramanujan-type identities, arXiv:2303.09069v3, 24 February 2025, 51 pp. https://arxiv.org/abs/2303.09069. Exact frozen author PDF reviewed; submission history consulted 27 October 2026.\par
EMS Press, Article record for doi:10.4171/JEMS/1627, https://ems.press/journals/jems/articles/14298710. Consulted 27 October 2026; publisher PDF access not established.\par
\end{document}
